\documentclass{article}
\usepackage{assignment_preamble}

\title{Homework 1}
\author{Ravi Kini}
\date{November 4, 2024}

\begin{document}

\maketitle

\problem
\subproblem{(a)}
Balanced bilingualism occurs when a bilingual has roughly equal proficiency in both of their languages, while unbalanced bilingualism occus when a bilingual has a stronger proficiency in one language in comparison to the other.
\subproblem{(b)}
Simultaneous bilingualism occurs when both languages are acquired simultaneously, while sequential bilingualism occurs when one language acquired after the other.
\subproblem{(c)}
Early bilingualism occurs when a child acquires both languages before the critical cut-off point while late bilingualism occurs when both languages are acquired after the critical cut-off point.

\clearpage

\problem
Bilingualism is referred to as a continuum because bilinguals vary in backgrounds, learning history, contexts of usage, and levels of competence.

\clearpage

\problem
Bilinguals tend to have slower processing speed, greater tolerance for ungrammaticality, and lower tolerance for ambiguity in comparison to monolinguals due to having to account for the syntax, lexicon, and disambiguation of both of their languages.

\clearpage

\problem
Since bilinguals differ from monolinguals in both of their languages in terms of knowledge, metalinguistic awareness, and cognitive processes, a bilingual cannot be considered as ``two monolinguals in one''.

\clearpage

\problem
Second language acquisition (SLA) is the late acquision of an L2 in either an L1 or L2 environment, usually involving education or immigration.

\clearpage

\problem
The critical period hypothesis is the claim that the period during which two languages can be acquired to monolingual levels of competence ends at a young age, estimated to be between two to eight years old.

\clearpage

\problem
Heritage bilingualism is the acquisition of both a home language and an environment language, and usually happens early, while second language acquisition is the acqusition of an L2 in an L1 or L2 environment, and usually happens later in life.

\clearpage

\problem
Typological factors can determine output in a second language by causing an ungrammatical word order. For example, speakers of Japanese, a SOV language, may place the verb at the end of a sentence when speaking English, a SVO language.

Internal factors can determine output in a second language by causing a speaker to have a higher level of competency in certain contexts. For example, early bilinguals tend to be more successful in informal and naturalistic contexts.

External factors can determine output in a second language by causing a speaker to adjust their speech to those with whom they converse. For example, according to communication accomodation theory, speakers may speak slower and more clearly and simplify their sentences in order to be better understood.

\clearpage

\problem
The CASP for Bilingualism principle of maximizing common ground refers to how bilinguals tend to maximize the common grammatical and lecial constructions of their two languages. For example, bilinguals speaking to other bilinguals that share the same languages are able to maximize the common ground between their two languages depending on their proficiency, while bilinguals speaking to monolinguals in one of their languages have less common ground available, and are more dependent on their proficiency.

\clearpage

\problem
An example of phonological transfer is when Spanish L1 speakers who are learning English add vowels before words that begin with an ``s'' and a consonant, since Spanish does not allow such sequences. An example of morphological transfer is English L2 speakers marking the past tense on both the main and auxiliary verb, something that is not done in English but is done in various other languages. An example of syntactic transfer is French L1 speakers who are learning English placng the verb before the adverb, as is done in French. An example of lexical transfer is Spanish L1 speakers who are learning English using false cognates like ``embarassed'' and ``embarasada''.

\clearpage

\problem
\subproblem{(a)}
\begin{exe}
    \ex[*] {I by car go to work.}
    \ex[*] {Today arrives my mum.}
\end{exe}
Since an error has been made in the word order, this error occured due to syntactic transfer. We can conclude that the L1 of the speaker is a Pro-Drop language.

\clearpage

\problem
U-shaped learning occurs when language learners initially produce correct forms by imitating native speakers, later produce incorrect forms as they learn the grammatical system, then producing the correct form again after the grammatical system has been cmpletely understood. For example, a child learning English may initally imitate the adults in their environment to produce sentences with irregular grammar patterns like ``many'' and ``more'', produce sentences with forms like ``many-er'' that follow expected grammatical patterns, then later revert to ``more'' as they come to understand which words are exceptions. This occurs because language acquisition does not always occur linearly, with learners gaining and losing skills.

\clearpage

\problem
More frequent forms are not always learnt sooner or better as a speaker may not always completely acquire a language.

\clearpage

\problem
Younger L2 learners have higher brain plasticity in comparison to older adult L2 learners, makign acquisition and storage of knowldge easier, but adult L2 learners can have greater analytical and pragmetic skills.

\clearpage

\problem
In bilingual production, restructuring occurs when two languages are incompatible, requiring that the grammatical system of one language be changed in order for a speaker to express themselved. For example, to express that ``this hotel forbids dogs'' in German or ``he ran across the street'' in French, it would be necessary to restructure the sentence.

\clearpage

\problem
Language typology is the field of linguistics that classifies languages according to their structural features in order to compare them. In the context of bilingualism, language typology is critical to understanding the differences between languages and how they lead to variation in competency among bilinguals and in comparison to monolinguals. For example, speakers of a SOV language may tend to make syntactic errors when speaking a SVO language due to a difference in the word orders.

\clearpage

\problem
Talmy's semantic typology indicates that some languages block the expression of certain pieces of information in certain places, requiring paraphrasing on the part of the speaker. In the context of bilingualism, bilinguals may use paraphrases due to typological restrictions on their expression in a certain language. For example, to indicate that one ran into the house in Spanish, one would have to explicitly indicate the manner after the clause, as it is optional, as opposed to being in the verb.

\clearpage

\problem
Since some languages make it easier or more difficult to express intentionality, information about intentionality can be lost when an event is decribed in a language that lacks this feature. For example, verb endings exist in Spanish that indicate whether an action was taken intentionally or unintentionally, but to express this in English would require restructuring the sentence.

\end{document}